\documentclass[11pt,letterpaper]{article}
\usepackage[technical-reference]{cornell-notes}
\usepackage{needspace}

\title{Clause 31: Input Output Library - Cornell Notes}
\author{}
\date{}
\setDocTitle{Clause 31: Input Output Library}
\setDocSubtitle{C++ 2024 Cornell Notes}
\setDocOwner{C++ 2024 Cornell Notes}
\setDocAuthor{}
\setDocDate{}
\setCornellCollection{C++ 2024 Cornell Notes}
\setCornellUnitType{Clause}
\setCornellUnitNumber{31}
\setCornellUnitTitle{Input Output Library}
\setCornellCompanionLabel{C++ Technical Reference}
\hypersetup{pdftitle={Clause 31: Input Output Library - Cornell Notes},pdfsubject={C++ technical study notes},pdfauthor={}}

\begin{document}
\maketitle
\begin{CornellReferenceOrientationBox}
\textbf{Purpose.} Defines stream state, formatting, buffering, stream buffers, string and span streams, file streams, synchronized output, filesystem paths and operations, and C-library file integration.
\par\textbf{Role.} The formatted and unformatted I/O model plus portable filesystem access.
\par\textbf{Principal dependencies.} Uses the library-wide rules of Clause 16 together with the applicable core-language concepts and supporting library clauses.
\par\textbf{Primary audience.} programmers, stream-buffer authors, and filesystem implementers.
\par\textbf{Major subject groups.} General; Iostreams requirements; Forward declarations; Standard iostream objects; Iostreams base classes; Stream buffers; Formatting and manipulators; String-based streams; Span-based streams; File-based streams; Synchronized output streams; File systems; and related subclauses.
\end{CornellReferenceOrientationBox}
\section{Learning Objectives}
\begin{itemize}[label=\textcolor{CornellReferenceTeal}{$\blacktriangleright$}]
\item Define the governing ideas of General and apply its stated contract at a call site.
\item Identify the governing ideas of Iostreams requirements and apply its stated contract at a call site.
\item Distinguish the governing ideas of Forward declarations and apply its stated contract at a call site.
\item Explain the governing ideas of Standard iostream objects and apply its stated contract at a call site.
\item Trace the governing ideas of Iostreams base classes and apply its stated contract at a call site.
\item Apply the governing ideas of Stream buffers and apply its stated contract at a call site.
\item Analyze the governing ideas of Formatting and manipulators and apply its stated contract at a call site.
\item Evaluate the governing ideas of String-based streams and apply its stated contract at a call site.
\item Diagnose the governing ideas of Span-based streams and apply its stated contract at a call site.
\item Compare the governing ideas of File-based streams and apply its stated contract at a call site.
\item Verify the governing ideas of Synchronized output streams and apply its stated contract at a call site.
\item Relate the governing ideas of File systems and apply its stated contract at a call site.
\item Classify the governing ideas of C library files and apply its stated contract at a call site.
\item Justify the governing ideas of Imbue limitations and apply its stated contract at a call site.
\item Audit the governing ideas of Types and apply its stated contract at a call site.
\item Summarize the governing ideas of Positioning type limitations and apply its stated contract at a call site.
\end{itemize}
\section{Normative Structure Map}
\small\begin{longtable}{@{}>{\RaggedRight\arraybackslash}p{0.13\textwidth}>{\RaggedRight\arraybackslash}p{0.27\textwidth}>{\RaggedRight\arraybackslash}p{0.19\textwidth}>{\RaggedRight\arraybackslash}p{0.31\textwidth}@{}}
\toprule\textbf{Subclause} & \textbf{Subject} & \textbf{Rule category} & \textbf{Key dependencies / entities}\\\midrule\endfirsthead
\toprule\textbf{Subclause} & \textbf{Subject} & \textbf{Rule category} & \textbf{Key dependencies / entities}\\\midrule\endhead
\bottomrule\endfoot
31.1 & General & library semantics & Uses the library-wide rules of Clause 16 together with the applicable core-language concepts and supporting library clauses.\\
31.2 & Iostreams requirements & requirements & Imbue limitations, Types, Positioning type limitations, Thread safety\\
31.3 & Forward declarations & library semantics & Header <iosfwd> synopsis, Overview\\
31.4 & Standard iostream objects & library semantics & Header <iostream> synopsis, Overview, Narrow stream objects, Wide stream objects\\
31.5 & Iostreams base classes & library semantics & Header <ios> synopsis, Class ios\_\allowbreak{}base, Class template fpos, Class template basic\_\allowbreak{}ios, and related subclauses\\
31.6 & Stream buffers & library semantics & Header <streambuf> synopsis, Stream buffer requirements, Class template basic\_\allowbreak{}streambuf\\
31.7 & Formatting and manipulators & library semantics & Header <istream> synopsis, Header <ostream> synopsis, Header <iomanip> synopsis, Header <print> synopsis, and related subclauses\\
31.8 & String-based streams & library semantics & Header <sstream> synopsis, Class template basic\_\allowbreak{}stringbuf, Class template basic\_\allowbreak{}istringstream, Class template basic\_\allowbreak{}ostringstream, and related subclauses\\
31.9 & Span-based streams & library semantics & Overview, Header <spanstream> synopsis, Class template basic\_\allowbreak{}spanbuf, Class template basic\_\allowbreak{}ispanstream, and related subclauses\\
31.10 & File-based streams & library semantics & Header <fstream> synopsis, Class template basic\_\allowbreak{}filebuf, Class template basic\_\allowbreak{}ifstream, Class template basic\_\allowbreak{}ofstream, and related subclauses\\
31.11 & Synchronized output streams & library semantics & Header <syncstream> synopsis, Class template basic\_\allowbreak{}syncbuf, Class template basic\_\allowbreak{}osyncstream\\
31.12 & File systems & library semantics & General, Conformance, Requirements, Header <filesystem> synopsis, and related subclauses\\
31.13 & C library files & library semantics & Header <cstdio> synopsis, Header <cinttypes> synopsis\\
\end{longtable}\normalsize
\begin{CornellReferenceNormativeBox}A heading maps the clause; it does not replace the detailed conditions. For each operation, read requirements, effects, result, exceptions, complexity, remarks, and notes with their stated normative force.\end{CornellReferenceNormativeBox}
\subsection*{Rule Classification Guide}
\small\begin{longtable}{@{}>{\RaggedRight\arraybackslash}p{0.25\textwidth}>{\RaggedRight\arraybackslash}p{0.67\textwidth}@{}}\toprule\textbf{Label or category} & \textbf{How to read it}\\\midrule\endfirsthead\toprule\textbf{Label or category} & \textbf{How to read it (continued)}\\\midrule\endhead\bottomrule\endfoot
Constraints & Conditions controlling whether a declaration, overload, or specialization participates in the relevant context.\\
Mandates & Requirements checked for the described declaration or specialization; failure has the stated well-formedness consequence.\\
Preconditions & Obligations the caller establishes before evaluation; violating one forfeits the operation's contract.\\
Effects / Returns / Postconditions & Separate the state change, produced result, and state guaranteed after successful completion.\\
Throws / Error conditions & State how failure is reported and which exceptional paths are permitted or required.\\
Complexity & Bounds or exact counts for the operations named by the specification.\\
Remarks / Recommended practice & Remarks refine applicability or participation; recommended practice guides implementations without silently becoming a program requirement.\\
Exposition-only & A device for describing semantics, not a public name on which a program can depend.\\
\end{longtable}\normalsize
\section{Cornell Cue-and-Notes}
\Needspace{8\baselineskip}
\subsection*{31.1 General}
\begin{longtable}{@{}>{\RaggedRight\arraybackslash\columncolor{CornellReferenceCueGray}}p{0.28\textwidth}>{\RaggedRight\arraybackslash}p{0.64\textwidth}@{}}
\rowcolor{CornellReferenceNavy}\color{white}\textbf{Cue / recall prompt} & \color{white}\textbf{Notes}\\\endfirsthead
\rowcolor{CornellReferenceNavy}\color{white}\textbf{Cue / recall prompt} & \color{white}\textbf{Notes (continued)}\\\endhead
\arrayrulecolor{CornellReferenceLineGray}
\textbf{What contract governs General?}\\[-1mm]{\scriptsize\CornellClauseRef{31.1}} & Frames the terminology, applicability, and relationships needed to interpret General; detailed rules remain in the following subclauses.\\\hline
\end{longtable}
\begin{CornellReferenceSummaryBox}General supplies the library semantics portion of this clause. The key review move is to connect its local wording to the clause-wide model, then classify each condition by normative force before relying on it.\end{CornellReferenceSummaryBox}
\Needspace{8\baselineskip}
\subsection*{31.2 Iostreams requirements}
\begin{longtable}{@{}>{\RaggedRight\arraybackslash\columncolor{CornellReferenceCueGray}}p{0.28\textwidth}>{\RaggedRight\arraybackslash}p{0.64\textwidth}@{}}
\rowcolor{CornellReferenceNavy}\color{white}\textbf{Cue / recall prompt} & \color{white}\textbf{Notes}\\\endfirsthead
\rowcolor{CornellReferenceNavy}\color{white}\textbf{Cue / recall prompt} & \color{white}\textbf{Notes (continued)}\\\endhead
\arrayrulecolor{CornellReferenceLineGray}
\textbf{What contract governs Imbue limitations?}\\[-1mm]{\scriptsize\CornellClauseRef{31.2.1}} & Specifies the facility family Imbue limitations. Read its declarations together with mandates, preconditions, effects, results, exceptions, complexity, and lifetime rules.\\\hline
\textbf{What contract governs Types?}\\[-1mm]{\scriptsize\CornellClauseRef{31.2.2}} & Specifies the facility family Types. Read its declarations together with mandates, preconditions, effects, results, exceptions, complexity, and lifetime rules.\\\hline
\textbf{What contract governs Positioning type limitations?}\\[-1mm]{\scriptsize\CornellClauseRef{31.2.3}} & Specifies the facility family Positioning type limitations. Read its declarations together with mandates, preconditions, effects, results, exceptions, complexity, and lifetime rules.\\\hline
\textbf{What contract governs Thread safety?}\\[-1mm]{\scriptsize\CornellClauseRef{31.2.4}} & Specifies the facility family Thread safety. Read its declarations together with mandates, preconditions, effects, results, exceptions, complexity, and lifetime rules.\\\hline
\end{longtable}
\begin{CornellReferenceSummaryBox}Iostreams requirements supplies the requirements portion of this clause. The key review move is to connect its local wording to the clause-wide model, then classify each condition by normative force before relying on it.\end{CornellReferenceSummaryBox}
\Needspace{8\baselineskip}
\subsection*{31.3 Forward declarations}
\begin{longtable}{@{}>{\RaggedRight\arraybackslash\columncolor{CornellReferenceCueGray}}p{0.28\textwidth}>{\RaggedRight\arraybackslash}p{0.64\textwidth}@{}}
\rowcolor{CornellReferenceNavy}\color{white}\textbf{Cue / recall prompt} & \color{white}\textbf{Notes}\\\endfirsthead
\rowcolor{CornellReferenceNavy}\color{white}\textbf{Cue / recall prompt} & \color{white}\textbf{Notes (continued)}\\\endhead
\arrayrulecolor{CornellReferenceLineGray}
\textbf{What contract governs Header <iosfwd> synopsis?}\\[-1mm]{\scriptsize\CornellClauseRef{31.3.1}} & Catalogs the declarations made available by Header <iosfwd> synopsis. The synopsis is an interface map; the detailed specifications supply constraints and behavior.\\\hline
\textbf{What contract governs Overview?}\\[-1mm]{\scriptsize\CornellClauseRef{31.3.2}} & Frames the terminology, applicability, and relationships needed to interpret Overview; detailed rules remain in the following subclauses.\\\hline
\end{longtable}
\begin{CornellReferenceSummaryBox}Forward declarations supplies the library semantics portion of this clause. The key review move is to connect its local wording to the clause-wide model, then classify each condition by normative force before relying on it.\end{CornellReferenceSummaryBox}
\Needspace{8\baselineskip}
\subsection*{31.4 Standard iostream objects}
\begin{longtable}{@{}>{\RaggedRight\arraybackslash\columncolor{CornellReferenceCueGray}}p{0.28\textwidth}>{\RaggedRight\arraybackslash}p{0.64\textwidth}@{}}
\rowcolor{CornellReferenceNavy}\color{white}\textbf{Cue / recall prompt} & \color{white}\textbf{Notes}\\\endfirsthead
\rowcolor{CornellReferenceNavy}\color{white}\textbf{Cue / recall prompt} & \color{white}\textbf{Notes (continued)}\\\endhead
\arrayrulecolor{CornellReferenceLineGray}
\textbf{What contract governs Header <iostream> synopsis?}\\[-1mm]{\scriptsize\CornellClauseRef{31.4.1}} & Catalogs the declarations made available by Header <iostream> synopsis. The synopsis is an interface map; the detailed specifications supply constraints and behavior.\\\hline
\textbf{What contract governs Overview?}\\[-1mm]{\scriptsize\CornellClauseRef{31.4.2}} & Frames the terminology, applicability, and relationships needed to interpret Overview; detailed rules remain in the following subclauses.\\\hline
\textbf{What contract governs Narrow stream objects?}\\[-1mm]{\scriptsize\CornellClauseRef{31.4.3}} & Specifies the facility family Narrow stream objects. Read its declarations together with mandates, preconditions, effects, results, exceptions, complexity, and lifetime rules.\\\hline
\textbf{What contract governs Wide stream objects?}\\[-1mm]{\scriptsize\CornellClauseRef{31.4.4}} & Specifies the facility family Wide stream objects. Read its declarations together with mandates, preconditions, effects, results, exceptions, complexity, and lifetime rules.\\\hline
\end{longtable}
\begin{CornellReferenceSummaryBox}Standard iostream objects supplies the library semantics portion of this clause. The key review move is to connect its local wording to the clause-wide model, then classify each condition by normative force before relying on it.\end{CornellReferenceSummaryBox}
\Needspace{8\baselineskip}
\subsection*{31.5 Iostreams base classes}
\begin{longtable}{@{}>{\RaggedRight\arraybackslash\columncolor{CornellReferenceCueGray}}p{0.28\textwidth}>{\RaggedRight\arraybackslash}p{0.64\textwidth}@{}}
\rowcolor{CornellReferenceNavy}\color{white}\textbf{Cue / recall prompt} & \color{white}\textbf{Notes}\\\endfirsthead
\rowcolor{CornellReferenceNavy}\color{white}\textbf{Cue / recall prompt} & \color{white}\textbf{Notes (continued)}\\\endhead
\arrayrulecolor{CornellReferenceLineGray}
\textbf{What contract governs Header <ios> synopsis?}\\[-1mm]{\scriptsize\CornellClauseRef{31.5.1}} & Catalogs the declarations made available by Header <ios> synopsis. The synopsis is an interface map; the detailed specifications supply constraints and behavior.\\\hline
\textbf{What contract governs Class ios\_\allowbreak{}base?}\\[-1mm]{\scriptsize\CornellClauseRef{31.5.2}} & Specifies the facility family Class ios\_\allowbreak{}base. Read its declarations together with mandates, preconditions, effects, results, exceptions, complexity, and lifetime rules.\\\hline
\textbf{What contract governs Class template fpos?}\\[-1mm]{\scriptsize\CornellClauseRef{31.5.3}} & Specifies the facility family Class template fpos. Read its declarations together with mandates, preconditions, effects, results, exceptions, complexity, and lifetime rules.\\\hline
\textbf{What contract governs Class template basic\_\allowbreak{}ios?}\\[-1mm]{\scriptsize\CornellClauseRef{31.5.4}} & Specifies the facility family Class template basic\_\allowbreak{}ios. Read its declarations together with mandates, preconditions, effects, results, exceptions, complexity, and lifetime rules.\\\hline
\textbf{What contract governs ios\_\allowbreak{}base manipulators?}\\[-1mm]{\scriptsize\CornellClauseRef{31.5.5}} & Specifies the facility family ios\_\allowbreak{}base manipulators. Read its declarations together with mandates, preconditions, effects, results, exceptions, complexity, and lifetime rules.\\\hline
\textbf{What contract governs Error reporting?}\\[-1mm]{\scriptsize\CornellClauseRef{31.5.6}} & Defines the failure model for Error reporting, including how an error is represented, reported, propagated, or converted into a diagnostic consequence.\\\hline
\end{longtable}
\begin{CornellReferenceSummaryBox}Iostreams base classes supplies the library semantics portion of this clause. The key review move is to connect its local wording to the clause-wide model, then classify each condition by normative force before relying on it.\end{CornellReferenceSummaryBox}
\Needspace{5\baselineskip}\paragraph{Detailed coverage ledger.}
\begin{description}[style=nextline,leftmargin=2.8cm,font=\normalfont\bfseries\color{CornellReferenceTeal}]
\item[31.5.2.1] \textbf{General.} Frames the terminology, applicability, and relationships needed to interpret General; detailed rules remain in the following subclauses.
\item[31.5.2.2] \textbf{Types.} Specifies the facility family Types. Read its declarations together with mandates, preconditions, effects, results, exceptions, complexity, and lifetime rules.
\item[31.5.2.2.1] \textbf{Class ios\_\allowbreak{}base::failure.} Specifies the facility family Class ios\_\allowbreak{}base::failure. Read its declarations together with mandates, preconditions, effects, results, exceptions, complexity, and lifetime rules.
\item[31.5.2.2.2] \textbf{Type ios\_\allowbreak{}base::fmtflags.} Specifies the facility family Type ios\_\allowbreak{}base::fmtflags. Read its declarations together with mandates, preconditions, effects, results, exceptions, complexity, and lifetime rules.
\item[31.5.2.2.3] \textbf{Type ios\_\allowbreak{}base::iostate.} Specifies the facility family Type ios\_\allowbreak{}base::iostate. Read its declarations together with mandates, preconditions, effects, results, exceptions, complexity, and lifetime rules.
\item[31.5.2.2.4] \textbf{Type ios\_\allowbreak{}base::openmode.} Specifies the facility family Type ios\_\allowbreak{}base::openmode. Read its declarations together with mandates, preconditions, effects, results, exceptions, complexity, and lifetime rules.
\item[31.5.2.2.5] \textbf{Type ios\_\allowbreak{}base::seekdir.} Specifies the facility family Type ios\_\allowbreak{}base::seekdir. Read its declarations together with mandates, preconditions, effects, results, exceptions, complexity, and lifetime rules.
\item[31.5.2.2.6] \textbf{Class ios\_\allowbreak{}base::Init.} Specifies the facility family Class ios\_\allowbreak{}base::Init. Read its declarations together with mandates, preconditions, effects, results, exceptions, complexity, and lifetime rules.
\item[31.5.2.3] \textbf{State functions.} Specifies the facility family State functions. Read its declarations together with mandates, preconditions, effects, results, exceptions, complexity, and lifetime rules.
\item[31.5.2.4] \textbf{Functions.} Specifies the facility family Functions. Read its declarations together with mandates, preconditions, effects, results, exceptions, complexity, and lifetime rules.
\item[31.5.2.5] \textbf{Static members.} Specifies the facility family Static members. Read its declarations together with mandates, preconditions, effects, results, exceptions, complexity, and lifetime rules.
\item[31.5.2.6] \textbf{Storage functions.} Specifies the facility family Storage functions. Read its declarations together with mandates, preconditions, effects, results, exceptions, complexity, and lifetime rules.
\item[31.5.2.7] \textbf{Callbacks.} Specifies the facility family Callbacks. Read its declarations together with mandates, preconditions, effects, results, exceptions, complexity, and lifetime rules.
\item[31.5.2.8] \textbf{Constructors and destructor.} Specifies how Constructors and destructor establishes object state, including participation constraints, initialization effects, and exceptional exits.
\item[31.5.3.1] \textbf{General.} Frames the terminology, applicability, and relationships needed to interpret General; detailed rules remain in the following subclauses.
\item[31.5.3.2] \textbf{Members.} Specifies the facility family Members. Read its declarations together with mandates, preconditions, effects, results, exceptions, complexity, and lifetime rules.
\item[31.5.3.3] \textbf{Requirements.} States the admissibility and semantic obligations for Requirements. Separate compile-time participation rules from caller preconditions and runtime effects.
\item[31.5.4.1] \textbf{Overview.} Frames the terminology, applicability, and relationships needed to interpret Overview; detailed rules remain in the following subclauses.
\item[31.5.4.2] \textbf{Constructors.} Specifies how Constructors establishes object state, including participation constraints, initialization effects, and exceptional exits.
\item[31.5.4.3] \textbf{Member functions.} Specifies the facility family Member functions. Read its declarations together with mandates, preconditions, effects, results, exceptions, complexity, and lifetime rules.
\item[31.5.4.4] \textbf{Flags functions.} Specifies the facility family Flags functions. Read its declarations together with mandates, preconditions, effects, results, exceptions, complexity, and lifetime rules.
\item[31.5.5.1] \textbf{fmtflags manipulators.} Specifies the facility family fmtflags manipulators. Read its declarations together with mandates, preconditions, effects, results, exceptions, complexity, and lifetime rules.
\item[31.5.5.2] \textbf{adjustfield manipulators.} Specifies the facility family adjustfield manipulators. Read its declarations together with mandates, preconditions, effects, results, exceptions, complexity, and lifetime rules.
\item[31.5.5.3] \textbf{basefield manipulators.} Specifies the facility family basefield manipulators. Read its declarations together with mandates, preconditions, effects, results, exceptions, complexity, and lifetime rules.
\item[31.5.5.4] \textbf{floatfield manipulators.} Specifies the facility family floatfield manipulators. Read its declarations together with mandates, preconditions, effects, results, exceptions, complexity, and lifetime rules.
\end{description}
\Needspace{8\baselineskip}
\subsection*{31.6 Stream buffers}
\begin{longtable}{@{}>{\RaggedRight\arraybackslash\columncolor{CornellReferenceCueGray}}p{0.28\textwidth}>{\RaggedRight\arraybackslash}p{0.64\textwidth}@{}}
\rowcolor{CornellReferenceNavy}\color{white}\textbf{Cue / recall prompt} & \color{white}\textbf{Notes}\\\endfirsthead
\rowcolor{CornellReferenceNavy}\color{white}\textbf{Cue / recall prompt} & \color{white}\textbf{Notes (continued)}\\\endhead
\arrayrulecolor{CornellReferenceLineGray}
\textbf{What contract governs Header <streambuf> synopsis?}\\[-1mm]{\scriptsize\CornellClauseRef{31.6.1}} & Catalogs the declarations made available by Header <streambuf> synopsis. The synopsis is an interface map; the detailed specifications supply constraints and behavior.\\\hline
\textbf{Which obligations govern Stream buffer requirements?}\\[-1mm]{\scriptsize\CornellClauseRef{31.6.2}} & States the admissibility and semantic obligations for Stream buffer requirements. Separate compile-time participation rules from caller preconditions and runtime effects.\\\hline
\textbf{What contract governs Class template basic\_\allowbreak{}streambuf?}\\[-1mm]{\scriptsize\CornellClauseRef{31.6.3}} & Specifies the facility family Class template basic\_\allowbreak{}streambuf. Read its declarations together with mandates, preconditions, effects, results, exceptions, complexity, and lifetime rules.\\\hline
\end{longtable}
\begin{CornellReferenceSummaryBox}Stream buffers supplies the library semantics portion of this clause. The key review move is to connect its local wording to the clause-wide model, then classify each condition by normative force before relying on it.\end{CornellReferenceSummaryBox}
\Needspace{5\baselineskip}\paragraph{Detailed coverage ledger.}
\begin{description}[style=nextline,leftmargin=2.8cm,font=\normalfont\bfseries\color{CornellReferenceTeal}]
\item[31.6.3.1] \textbf{General.} Frames the terminology, applicability, and relationships needed to interpret General; detailed rules remain in the following subclauses.
\item[31.6.3.2] \textbf{Constructors.} Specifies how Constructors establishes object state, including participation constraints, initialization effects, and exceptional exits.
\item[31.6.3.3] \textbf{Public member functions.} Specifies the facility family Public member functions. Read its declarations together with mandates, preconditions, effects, results, exceptions, complexity, and lifetime rules.
\item[31.6.3.3.1] \textbf{Locales.} Specifies the facility family Locales. Read its declarations together with mandates, preconditions, effects, results, exceptions, complexity, and lifetime rules.
\item[31.6.3.3.2] \textbf{Buffer management and positioning.} Specifies the facility family Buffer management and positioning. Read its declarations together with mandates, preconditions, effects, results, exceptions, complexity, and lifetime rules.
\item[31.6.3.3.3] \textbf{Get area.} Specifies the facility family Get area. Read its declarations together with mandates, preconditions, effects, results, exceptions, complexity, and lifetime rules.
\item[31.6.3.3.4] \textbf{Putback.} Specifies the facility family Putback. Read its declarations together with mandates, preconditions, effects, results, exceptions, complexity, and lifetime rules.
\item[31.6.3.3.5] \textbf{Put area.} Specifies the facility family Put area. Read its declarations together with mandates, preconditions, effects, results, exceptions, complexity, and lifetime rules.
\item[31.6.3.4] \textbf{Protected member functions.} Specifies the facility family Protected member functions. Read its declarations together with mandates, preconditions, effects, results, exceptions, complexity, and lifetime rules.
\item[31.6.3.4.1] \textbf{Assignment.} Governs state transfer or exchange for Assignment; check self-operation, validity after movement, exception behavior, and invalidation.
\item[31.6.3.4.2] \textbf{Get area access.} Specifies the facility family Get area access. Read its declarations together with mandates, preconditions, effects, results, exceptions, complexity, and lifetime rules.
\item[31.6.3.4.3] \textbf{Put area access.} Specifies the facility family Put area access. Read its declarations together with mandates, preconditions, effects, results, exceptions, complexity, and lifetime rules.
\item[31.6.3.5] \textbf{Virtual functions.} Specifies the facility family Virtual functions. Read its declarations together with mandates, preconditions, effects, results, exceptions, complexity, and lifetime rules.
\item[31.6.3.5.1] \textbf{Locales.} Specifies the facility family Locales. Read its declarations together with mandates, preconditions, effects, results, exceptions, complexity, and lifetime rules.
\item[31.6.3.5.2] \textbf{Buffer management and positioning.} Specifies the facility family Buffer management and positioning. Read its declarations together with mandates, preconditions, effects, results, exceptions, complexity, and lifetime rules.
\item[31.6.3.5.3] \textbf{Get area.} Specifies the facility family Get area. Read its declarations together with mandates, preconditions, effects, results, exceptions, complexity, and lifetime rules.
\item[31.6.3.5.4] \textbf{Putback.} Specifies the facility family Putback. Read its declarations together with mandates, preconditions, effects, results, exceptions, complexity, and lifetime rules.
\item[31.6.3.5.5] \textbf{Put area.} Specifies the facility family Put area. Read its declarations together with mandates, preconditions, effects, results, exceptions, complexity, and lifetime rules.
\end{description}
\Needspace{8\baselineskip}
\subsection*{31.7 Formatting and manipulators}
\begin{longtable}{@{}>{\RaggedRight\arraybackslash\columncolor{CornellReferenceCueGray}}p{0.28\textwidth}>{\RaggedRight\arraybackslash}p{0.64\textwidth}@{}}
\rowcolor{CornellReferenceNavy}\color{white}\textbf{Cue / recall prompt} & \color{white}\textbf{Notes}\\\endfirsthead
\rowcolor{CornellReferenceNavy}\color{white}\textbf{Cue / recall prompt} & \color{white}\textbf{Notes (continued)}\\\endhead
\arrayrulecolor{CornellReferenceLineGray}
\textbf{What contract governs Header <istream> synopsis?}\\[-1mm]{\scriptsize\CornellClauseRef{31.7.1}} & Catalogs the declarations made available by Header <istream> synopsis. The synopsis is an interface map; the detailed specifications supply constraints and behavior.\\\hline
\textbf{What contract governs Header <ostream> synopsis?}\\[-1mm]{\scriptsize\CornellClauseRef{31.7.2}} & Catalogs the declarations made available by Header <ostream> synopsis. The synopsis is an interface map; the detailed specifications supply constraints and behavior.\\\hline
\textbf{What contract governs Header <iomanip> synopsis?}\\[-1mm]{\scriptsize\CornellClauseRef{31.7.3}} & Catalogs the declarations made available by Header <iomanip> synopsis. The synopsis is an interface map; the detailed specifications supply constraints and behavior.\\\hline
\textbf{What contract governs Header <print> synopsis?}\\[-1mm]{\scriptsize\CornellClauseRef{31.7.4}} & Catalogs the declarations made available by Header <print> synopsis. The synopsis is an interface map; the detailed specifications supply constraints and behavior.\\\hline
\textbf{What contract governs Input streams?}\\[-1mm]{\scriptsize\CornellClauseRef{31.7.5}} & Specifies the facility family Input streams. Read its declarations together with mandates, preconditions, effects, results, exceptions, complexity, and lifetime rules.\\\hline
\textbf{What contract governs Output streams?}\\[-1mm]{\scriptsize\CornellClauseRef{31.7.6}} & Specifies the facility family Output streams. Read its declarations together with mandates, preconditions, effects, results, exceptions, complexity, and lifetime rules.\\\hline
\textbf{What contract governs Standard manipulators?}\\[-1mm]{\scriptsize\CornellClauseRef{31.7.7}} & Specifies the facility family Standard manipulators. Read its declarations together with mandates, preconditions, effects, results, exceptions, complexity, and lifetime rules.\\\hline
\textbf{What contract governs Extended manipulators?}\\[-1mm]{\scriptsize\CornellClauseRef{31.7.8}} & Specifies the facility family Extended manipulators. Read its declarations together with mandates, preconditions, effects, results, exceptions, complexity, and lifetime rules.\\\hline
\textbf{What contract governs Quoted manipulators?}\\[-1mm]{\scriptsize\CornellClauseRef{31.7.9}} & Specifies the facility family Quoted manipulators. Read its declarations together with mandates, preconditions, effects, results, exceptions, complexity, and lifetime rules.\\\hline
\textbf{What contract governs Print functions?}\\[-1mm]{\scriptsize\CornellClauseRef{31.7.10}} & Specifies the facility family Print functions. Read its declarations together with mandates, preconditions, effects, results, exceptions, complexity, and lifetime rules.\\\hline
\end{longtable}
\begin{CornellReferenceSummaryBox}Formatting and manipulators supplies the library semantics portion of this clause. The key review move is to connect its local wording to the clause-wide model, then classify each condition by normative force before relying on it.\end{CornellReferenceSummaryBox}
\Needspace{5\baselineskip}\paragraph{Detailed coverage ledger.}
\begin{description}[style=nextline,leftmargin=2.8cm,font=\normalfont\bfseries\color{CornellReferenceTeal}]
\item[31.7.5.1] \textbf{General.} Frames the terminology, applicability, and relationships needed to interpret General; detailed rules remain in the following subclauses.
\item[31.7.5.2] \textbf{Class template basic\_\allowbreak{}istream.} Specifies the facility family Class template basic\_\allowbreak{}istream. Read its declarations together with mandates, preconditions, effects, results, exceptions, complexity, and lifetime rules.
\item[31.7.5.2.1] \textbf{General.} Frames the terminology, applicability, and relationships needed to interpret General; detailed rules remain in the following subclauses.
\item[31.7.5.2.2] \textbf{Constructors.} Specifies how Constructors establishes object state, including participation constraints, initialization effects, and exceptional exits.
\item[31.7.5.2.3] \textbf{Assignment and swap.} Governs state transfer or exchange for Assignment and swap; check self-operation, validity after movement, exception behavior, and invalidation.
\item[31.7.5.2.4] \textbf{Class basic\_\allowbreak{}istream::sentry.} Specifies the facility family Class basic\_\allowbreak{}istream::sentry. Read its declarations together with mandates, preconditions, effects, results, exceptions, complexity, and lifetime rules.
\item[31.7.5.3] \textbf{Formatted input functions.} Specifies the facility family Formatted input functions. Read its declarations together with mandates, preconditions, effects, results, exceptions, complexity, and lifetime rules.
\item[31.7.5.3.1] \textbf{Common requirements.} States the admissibility and semantic obligations for Common requirements. Separate compile-time participation rules from caller preconditions and runtime effects.
\item[31.7.5.3.2] \textbf{Arithmetic extractors.} Specifies the facility family Arithmetic extractors. Read its declarations together with mandates, preconditions, effects, results, exceptions, complexity, and lifetime rules.
\item[31.7.5.3.3] \textbf{basic\_\allowbreak{}istream::operator>>.} Specifies the facility family basic\_\allowbreak{}istream::operator>>. Read its declarations together with mandates, preconditions, effects, results, exceptions, complexity, and lifetime rules.
\item[31.7.5.4] \textbf{Unformatted input functions.} Specifies the facility family Unformatted input functions. Read its declarations together with mandates, preconditions, effects, results, exceptions, complexity, and lifetime rules.
\item[31.7.5.5] \textbf{Standard basic\_\allowbreak{}istream manipulators.} Specifies the facility family Standard basic\_\allowbreak{}istream manipulators. Read its declarations together with mandates, preconditions, effects, results, exceptions, complexity, and lifetime rules.
\item[31.7.5.6] \textbf{Rvalue stream extraction.} Specifies the facility family Rvalue stream extraction. Read its declarations together with mandates, preconditions, effects, results, exceptions, complexity, and lifetime rules.
\item[31.7.5.7] \textbf{Class template basic\_\allowbreak{}iostream.} Specifies the facility family Class template basic\_\allowbreak{}iostream. Read its declarations together with mandates, preconditions, effects, results, exceptions, complexity, and lifetime rules.
\item[31.7.5.7.1] \textbf{General.} Frames the terminology, applicability, and relationships needed to interpret General; detailed rules remain in the following subclauses.
\item[31.7.5.7.2] \textbf{Constructors.} Specifies how Constructors establishes object state, including participation constraints, initialization effects, and exceptional exits.
\item[31.7.5.7.3] \textbf{Destructor.} Specifies how Destructor ends the object's managed state and which cleanup or termination consequences apply.
\item[31.7.5.7.4] \textbf{Assignment and swap.} Governs state transfer or exchange for Assignment and swap; check self-operation, validity after movement, exception behavior, and invalidation.
\item[31.7.6.1] \textbf{General.} Frames the terminology, applicability, and relationships needed to interpret General; detailed rules remain in the following subclauses.
\item[31.7.6.2] \textbf{Class template basic\_\allowbreak{}ostream.} Specifies the facility family Class template basic\_\allowbreak{}ostream. Read its declarations together with mandates, preconditions, effects, results, exceptions, complexity, and lifetime rules.
\item[31.7.6.2.1] \textbf{General.} Frames the terminology, applicability, and relationships needed to interpret General; detailed rules remain in the following subclauses.
\item[31.7.6.2.2] \textbf{Constructors.} Specifies how Constructors establishes object state, including participation constraints, initialization effects, and exceptional exits.
\item[31.7.6.2.3] \textbf{Assignment and swap.} Governs state transfer or exchange for Assignment and swap; check self-operation, validity after movement, exception behavior, and invalidation.
\item[31.7.6.2.4] \textbf{Class basic\_\allowbreak{}ostream::sentry.} Specifies the facility family Class basic\_\allowbreak{}ostream::sentry. Read its declarations together with mandates, preconditions, effects, results, exceptions, complexity, and lifetime rules.
\item[31.7.6.2.5] \textbf{Seek members.} Specifies the facility family Seek members. Read its declarations together with mandates, preconditions, effects, results, exceptions, complexity, and lifetime rules.
\item[31.7.6.3] \textbf{Formatted output functions.} Specifies the facility family Formatted output functions. Read its declarations together with mandates, preconditions, effects, results, exceptions, complexity, and lifetime rules.
\item[31.7.6.3.1] \textbf{Common requirements.} States the admissibility and semantic obligations for Common requirements. Separate compile-time participation rules from caller preconditions and runtime effects.
\item[31.7.6.3.2] \textbf{Arithmetic inserters.} Specifies the facility family Arithmetic inserters. Read its declarations together with mandates, preconditions, effects, results, exceptions, complexity, and lifetime rules.
\item[31.7.6.3.3] \textbf{basic\_\allowbreak{}ostream::operator<<.} Specifies the facility family basic\_\allowbreak{}ostream::operator<<. Read its declarations together with mandates, preconditions, effects, results, exceptions, complexity, and lifetime rules.
\item[31.7.6.3.4] \textbf{Character inserter function templates.} Specifies text-sequence behavior for Character inserter function templates; establish character type, extent, ownership, null termination, encoding assumptions, and invalidation before use.
\item[31.7.6.3.5] \textbf{Print.} Specifies the facility family Print. Read its declarations together with mandates, preconditions, effects, results, exceptions, complexity, and lifetime rules.
\item[31.7.6.4] \textbf{Unformatted output functions.} Specifies the facility family Unformatted output functions. Read its declarations together with mandates, preconditions, effects, results, exceptions, complexity, and lifetime rules.
\item[31.7.6.5] \textbf{Standard manipulators.} Specifies the facility family Standard manipulators. Read its declarations together with mandates, preconditions, effects, results, exceptions, complexity, and lifetime rules.
\item[31.7.6.6] \textbf{Rvalue stream insertion.} Specifies the facility family Rvalue stream insertion. Read its declarations together with mandates, preconditions, effects, results, exceptions, complexity, and lifetime rules.
\end{description}
\Needspace{8\baselineskip}
\subsection*{31.8 String-based streams}
\begin{longtable}{@{}>{\RaggedRight\arraybackslash\columncolor{CornellReferenceCueGray}}p{0.28\textwidth}>{\RaggedRight\arraybackslash}p{0.64\textwidth}@{}}
\rowcolor{CornellReferenceNavy}\color{white}\textbf{Cue / recall prompt} & \color{white}\textbf{Notes}\\\endfirsthead
\rowcolor{CornellReferenceNavy}\color{white}\textbf{Cue / recall prompt} & \color{white}\textbf{Notes (continued)}\\\endhead
\arrayrulecolor{CornellReferenceLineGray}
\textbf{What contract governs Header <sstream> synopsis?}\\[-1mm]{\scriptsize\CornellClauseRef{31.8.1}} & Catalogs the declarations made available by Header <sstream> synopsis. The synopsis is an interface map; the detailed specifications supply constraints and behavior.\\\hline
\textbf{What contract governs Class template basic\_\allowbreak{}stringbuf?}\\[-1mm]{\scriptsize\CornellClauseRef{31.8.2}} & Specifies text-sequence behavior for Class template basic\_\allowbreak{}stringbuf; establish character type, extent, ownership, null termination, encoding assumptions, and invalidation before use.\\\hline
\textbf{What contract governs Class template basic\_\allowbreak{}istringstream?}\\[-1mm]{\scriptsize\CornellClauseRef{31.8.3}} & Specifies text-sequence behavior for Class template basic\_\allowbreak{}istringstream; establish character type, extent, ownership, null termination, encoding assumptions, and invalidation before use.\\\hline
\textbf{What contract governs Class template basic\_\allowbreak{}ostringstream?}\\[-1mm]{\scriptsize\CornellClauseRef{31.8.4}} & Specifies text-sequence behavior for Class template basic\_\allowbreak{}ostringstream; establish character type, extent, ownership, null termination, encoding assumptions, and invalidation before use.\\\hline
\textbf{What contract governs Class template basic\_\allowbreak{}stringstream?}\\[-1mm]{\scriptsize\CornellClauseRef{31.8.5}} & Specifies text-sequence behavior for Class template basic\_\allowbreak{}stringstream; establish character type, extent, ownership, null termination, encoding assumptions, and invalidation before use.\\\hline
\end{longtable}
\begin{CornellReferenceSummaryBox}String-based streams supplies the library semantics portion of this clause. The key review move is to connect its local wording to the clause-wide model, then classify each condition by normative force before relying on it.\end{CornellReferenceSummaryBox}
\Needspace{5\baselineskip}\paragraph{Detailed coverage ledger.}
\begin{description}[style=nextline,leftmargin=2.8cm,font=\normalfont\bfseries\color{CornellReferenceTeal}]
\item[31.8.2.1] \textbf{General.} Frames the terminology, applicability, and relationships needed to interpret General; detailed rules remain in the following subclauses.
\item[31.8.2.2] \textbf{Constructors.} Specifies how Constructors establishes object state, including participation constraints, initialization effects, and exceptional exits.
\item[31.8.2.3] \textbf{Assignment and swap.} Governs state transfer or exchange for Assignment and swap; check self-operation, validity after movement, exception behavior, and invalidation.
\item[31.8.2.4] \textbf{Member functions.} Specifies the facility family Member functions. Read its declarations together with mandates, preconditions, effects, results, exceptions, complexity, and lifetime rules.
\item[31.8.2.5] \textbf{Overridden virtual functions.} Specifies the facility family Overridden virtual functions. Read its declarations together with mandates, preconditions, effects, results, exceptions, complexity, and lifetime rules.
\item[31.8.3.1] \textbf{General.} Frames the terminology, applicability, and relationships needed to interpret General; detailed rules remain in the following subclauses.
\item[31.8.3.2] \textbf{Constructors.} Specifies how Constructors establishes object state, including participation constraints, initialization effects, and exceptional exits.
\item[31.8.3.3] \textbf{Swap.} Governs state transfer or exchange for Swap; check self-operation, validity after movement, exception behavior, and invalidation.
\item[31.8.3.4] \textbf{Member functions.} Specifies the facility family Member functions. Read its declarations together with mandates, preconditions, effects, results, exceptions, complexity, and lifetime rules.
\item[31.8.4.1] \textbf{General.} Frames the terminology, applicability, and relationships needed to interpret General; detailed rules remain in the following subclauses.
\item[31.8.4.2] \textbf{Constructors.} Specifies how Constructors establishes object state, including participation constraints, initialization effects, and exceptional exits.
\item[31.8.4.3] \textbf{Swap.} Governs state transfer or exchange for Swap; check self-operation, validity after movement, exception behavior, and invalidation.
\item[31.8.4.4] \textbf{Member functions.} Specifies the facility family Member functions. Read its declarations together with mandates, preconditions, effects, results, exceptions, complexity, and lifetime rules.
\item[31.8.5.1] \textbf{General.} Frames the terminology, applicability, and relationships needed to interpret General; detailed rules remain in the following subclauses.
\item[31.8.5.2] \textbf{Constructors.} Specifies how Constructors establishes object state, including participation constraints, initialization effects, and exceptional exits.
\item[31.8.5.3] \textbf{Swap.} Governs state transfer or exchange for Swap; check self-operation, validity after movement, exception behavior, and invalidation.
\item[31.8.5.4] \textbf{Member functions.} Specifies the facility family Member functions. Read its declarations together with mandates, preconditions, effects, results, exceptions, complexity, and lifetime rules.
\end{description}
\Needspace{8\baselineskip}
\subsection*{31.9 Span-based streams}
\begin{longtable}{@{}>{\RaggedRight\arraybackslash\columncolor{CornellReferenceCueGray}}p{0.28\textwidth}>{\RaggedRight\arraybackslash}p{0.64\textwidth}@{}}
\rowcolor{CornellReferenceNavy}\color{white}\textbf{Cue / recall prompt} & \color{white}\textbf{Notes}\\\endfirsthead
\rowcolor{CornellReferenceNavy}\color{white}\textbf{Cue / recall prompt} & \color{white}\textbf{Notes (continued)}\\\endhead
\arrayrulecolor{CornellReferenceLineGray}
\textbf{What contract governs Overview?}\\[-1mm]{\scriptsize\CornellClauseRef{31.9.1}} & Frames the terminology, applicability, and relationships needed to interpret Overview; detailed rules remain in the following subclauses.\\\hline
\textbf{What contract governs Header <spanstream> synopsis?}\\[-1mm]{\scriptsize\CornellClauseRef{31.9.2}} & Catalogs the declarations made available by Header <spanstream> synopsis. The synopsis is an interface map; the detailed specifications supply constraints and behavior.\\\hline
\textbf{What contract governs Class template basic\_\allowbreak{}spanbuf?}\\[-1mm]{\scriptsize\CornellClauseRef{31.9.3}} & Specifies the facility family Class template basic\_\allowbreak{}spanbuf. Read its declarations together with mandates, preconditions, effects, results, exceptions, complexity, and lifetime rules.\\\hline
\textbf{What contract governs Class template basic\_\allowbreak{}ispanstream?}\\[-1mm]{\scriptsize\CornellClauseRef{31.9.4}} & Specifies the facility family Class template basic\_\allowbreak{}ispanstream. Read its declarations together with mandates, preconditions, effects, results, exceptions, complexity, and lifetime rules.\\\hline
\textbf{What contract governs Class template basic\_\allowbreak{}ospanstream?}\\[-1mm]{\scriptsize\CornellClauseRef{31.9.5}} & Specifies the facility family Class template basic\_\allowbreak{}ospanstream. Read its declarations together with mandates, preconditions, effects, results, exceptions, complexity, and lifetime rules.\\\hline
\textbf{What contract governs Class template basic\_\allowbreak{}spanstream?}\\[-1mm]{\scriptsize\CornellClauseRef{31.9.6}} & Specifies the facility family Class template basic\_\allowbreak{}spanstream. Read its declarations together with mandates, preconditions, effects, results, exceptions, complexity, and lifetime rules.\\\hline
\end{longtable}
\begin{CornellReferenceSummaryBox}Span-based streams supplies the library semantics portion of this clause. The key review move is to connect its local wording to the clause-wide model, then classify each condition by normative force before relying on it.\end{CornellReferenceSummaryBox}
\Needspace{5\baselineskip}\paragraph{Detailed coverage ledger.}
\begin{description}[style=nextline,leftmargin=2.8cm,font=\normalfont\bfseries\color{CornellReferenceTeal}]
\item[31.9.3.1] \textbf{General.} Frames the terminology, applicability, and relationships needed to interpret General; detailed rules remain in the following subclauses.
\item[31.9.3.2] \textbf{Constructors.} Specifies how Constructors establishes object state, including participation constraints, initialization effects, and exceptional exits.
\item[31.9.3.3] \textbf{Assignment and swap.} Governs state transfer or exchange for Assignment and swap; check self-operation, validity after movement, exception behavior, and invalidation.
\item[31.9.3.4] \textbf{Member functions.} Specifies the facility family Member functions. Read its declarations together with mandates, preconditions, effects, results, exceptions, complexity, and lifetime rules.
\item[31.9.3.5] \textbf{Overridden virtual functions.} Specifies the facility family Overridden virtual functions. Read its declarations together with mandates, preconditions, effects, results, exceptions, complexity, and lifetime rules.
\item[31.9.4.1] \textbf{General.} Frames the terminology, applicability, and relationships needed to interpret General; detailed rules remain in the following subclauses.
\item[31.9.4.2] \textbf{Constructors.} Specifies how Constructors establishes object state, including participation constraints, initialization effects, and exceptional exits.
\item[31.9.4.3] \textbf{Swap.} Governs state transfer or exchange for Swap; check self-operation, validity after movement, exception behavior, and invalidation.
\item[31.9.4.4] \textbf{Member functions.} Specifies the facility family Member functions. Read its declarations together with mandates, preconditions, effects, results, exceptions, complexity, and lifetime rules.
\item[31.9.5.1] \textbf{General.} Frames the terminology, applicability, and relationships needed to interpret General; detailed rules remain in the following subclauses.
\item[31.9.5.2] \textbf{Constructors.} Specifies how Constructors establishes object state, including participation constraints, initialization effects, and exceptional exits.
\item[31.9.5.3] \textbf{Swap.} Governs state transfer or exchange for Swap; check self-operation, validity after movement, exception behavior, and invalidation.
\item[31.9.5.4] \textbf{Member functions.} Specifies the facility family Member functions. Read its declarations together with mandates, preconditions, effects, results, exceptions, complexity, and lifetime rules.
\item[31.9.6.1] \textbf{General.} Frames the terminology, applicability, and relationships needed to interpret General; detailed rules remain in the following subclauses.
\item[31.9.6.2] \textbf{Constructors.} Specifies how Constructors establishes object state, including participation constraints, initialization effects, and exceptional exits.
\item[31.9.6.3] \textbf{Swap.} Governs state transfer or exchange for Swap; check self-operation, validity after movement, exception behavior, and invalidation.
\item[31.9.6.4] \textbf{Member functions.} Specifies the facility family Member functions. Read its declarations together with mandates, preconditions, effects, results, exceptions, complexity, and lifetime rules.
\end{description}
\Needspace{8\baselineskip}
\subsection*{31.10 File-based streams}
\begin{longtable}{@{}>{\RaggedRight\arraybackslash\columncolor{CornellReferenceCueGray}}p{0.28\textwidth}>{\RaggedRight\arraybackslash}p{0.64\textwidth}@{}}
\rowcolor{CornellReferenceNavy}\color{white}\textbf{Cue / recall prompt} & \color{white}\textbf{Notes}\\\endfirsthead
\rowcolor{CornellReferenceNavy}\color{white}\textbf{Cue / recall prompt} & \color{white}\textbf{Notes (continued)}\\\endhead
\arrayrulecolor{CornellReferenceLineGray}
\textbf{What contract governs Header <fstream> synopsis?}\\[-1mm]{\scriptsize\CornellClauseRef{31.10.1}} & Catalogs the declarations made available by Header <fstream> synopsis. The synopsis is an interface map; the detailed specifications supply constraints and behavior.\\\hline
\textbf{What contract governs Class template basic\_\allowbreak{}filebuf?}\\[-1mm]{\scriptsize\CornellClauseRef{31.10.2}} & Specifies the facility family Class template basic\_\allowbreak{}filebuf. Read its declarations together with mandates, preconditions, effects, results, exceptions, complexity, and lifetime rules.\\\hline
\textbf{What contract governs Class template basic\_\allowbreak{}ifstream?}\\[-1mm]{\scriptsize\CornellClauseRef{31.10.3}} & Specifies the facility family Class template basic\_\allowbreak{}ifstream. Read its declarations together with mandates, preconditions, effects, results, exceptions, complexity, and lifetime rules.\\\hline
\textbf{What contract governs Class template basic\_\allowbreak{}ofstream?}\\[-1mm]{\scriptsize\CornellClauseRef{31.10.4}} & Specifies the facility family Class template basic\_\allowbreak{}ofstream. Read its declarations together with mandates, preconditions, effects, results, exceptions, complexity, and lifetime rules.\\\hline
\textbf{What contract governs Class template basic\_\allowbreak{}fstream?}\\[-1mm]{\scriptsize\CornellClauseRef{31.10.5}} & Specifies the facility family Class template basic\_\allowbreak{}fstream. Read its declarations together with mandates, preconditions, effects, results, exceptions, complexity, and lifetime rules.\\\hline
\end{longtable}
\begin{CornellReferenceSummaryBox}File-based streams supplies the library semantics portion of this clause. The key review move is to connect its local wording to the clause-wide model, then classify each condition by normative force before relying on it.\end{CornellReferenceSummaryBox}
\Needspace{5\baselineskip}\paragraph{Detailed coverage ledger.}
\begin{description}[style=nextline,leftmargin=2.8cm,font=\normalfont\bfseries\color{CornellReferenceTeal}]
\item[31.10.2.1] \textbf{General.} Frames the terminology, applicability, and relationships needed to interpret General; detailed rules remain in the following subclauses.
\item[31.10.2.2] \textbf{Constructors.} Specifies how Constructors establishes object state, including participation constraints, initialization effects, and exceptional exits.
\item[31.10.2.3] \textbf{Assignment and swap.} Governs state transfer or exchange for Assignment and swap; check self-operation, validity after movement, exception behavior, and invalidation.
\item[31.10.2.4] \textbf{Member functions.} Specifies the facility family Member functions. Read its declarations together with mandates, preconditions, effects, results, exceptions, complexity, and lifetime rules.
\item[31.10.2.5] \textbf{Overridden virtual functions.} Specifies the facility family Overridden virtual functions. Read its declarations together with mandates, preconditions, effects, results, exceptions, complexity, and lifetime rules.
\item[31.10.3.1] \textbf{General.} Frames the terminology, applicability, and relationships needed to interpret General; detailed rules remain in the following subclauses.
\item[31.10.3.2] \textbf{Constructors.} Specifies how Constructors establishes object state, including participation constraints, initialization effects, and exceptional exits.
\item[31.10.3.3] \textbf{Swap.} Governs state transfer or exchange for Swap; check self-operation, validity after movement, exception behavior, and invalidation.
\item[31.10.3.4] \textbf{Member functions.} Specifies the facility family Member functions. Read its declarations together with mandates, preconditions, effects, results, exceptions, complexity, and lifetime rules.
\item[31.10.4.1] \textbf{General.} Frames the terminology, applicability, and relationships needed to interpret General; detailed rules remain in the following subclauses.
\item[31.10.4.2] \textbf{Constructors.} Specifies how Constructors establishes object state, including participation constraints, initialization effects, and exceptional exits.
\item[31.10.4.3] \textbf{Swap.} Governs state transfer or exchange for Swap; check self-operation, validity after movement, exception behavior, and invalidation.
\item[31.10.4.4] \textbf{Member functions.} Specifies the facility family Member functions. Read its declarations together with mandates, preconditions, effects, results, exceptions, complexity, and lifetime rules.
\item[31.10.5.1] \textbf{General.} Frames the terminology, applicability, and relationships needed to interpret General; detailed rules remain in the following subclauses.
\item[31.10.5.2] \textbf{Constructors.} Specifies how Constructors establishes object state, including participation constraints, initialization effects, and exceptional exits.
\item[31.10.5.3] \textbf{Swap.} Governs state transfer or exchange for Swap; check self-operation, validity after movement, exception behavior, and invalidation.
\item[31.10.5.4] \textbf{Member functions.} Specifies the facility family Member functions. Read its declarations together with mandates, preconditions, effects, results, exceptions, complexity, and lifetime rules.
\end{description}
\Needspace{8\baselineskip}
\subsection*{31.11 Synchronized output streams}
\begin{longtable}{@{}>{\RaggedRight\arraybackslash\columncolor{CornellReferenceCueGray}}p{0.28\textwidth}>{\RaggedRight\arraybackslash}p{0.64\textwidth}@{}}
\rowcolor{CornellReferenceNavy}\color{white}\textbf{Cue / recall prompt} & \color{white}\textbf{Notes}\\\endfirsthead
\rowcolor{CornellReferenceNavy}\color{white}\textbf{Cue / recall prompt} & \color{white}\textbf{Notes (continued)}\\\endhead
\arrayrulecolor{CornellReferenceLineGray}
\textbf{What contract governs Header <syncstream> synopsis?}\\[-1mm]{\scriptsize\CornellClauseRef{31.11.1}} & Catalogs the declarations made available by Header <syncstream> synopsis. The synopsis is an interface map; the detailed specifications supply constraints and behavior.\\\hline
\textbf{What contract governs Class template basic\_\allowbreak{}syncbuf?}\\[-1mm]{\scriptsize\CornellClauseRef{31.11.2}} & Specifies the facility family Class template basic\_\allowbreak{}syncbuf. Read its declarations together with mandates, preconditions, effects, results, exceptions, complexity, and lifetime rules.\\\hline
\textbf{What contract governs Class template basic\_\allowbreak{}osyncstream?}\\[-1mm]{\scriptsize\CornellClauseRef{31.11.3}} & Specifies the facility family Class template basic\_\allowbreak{}osyncstream. Read its declarations together with mandates, preconditions, effects, results, exceptions, complexity, and lifetime rules.\\\hline
\end{longtable}
\begin{CornellReferenceSummaryBox}Synchronized output streams supplies the library semantics portion of this clause. The key review move is to connect its local wording to the clause-wide model, then classify each condition by normative force before relying on it.\end{CornellReferenceSummaryBox}
\Needspace{5\baselineskip}\paragraph{Detailed coverage ledger.}
\begin{description}[style=nextline,leftmargin=2.8cm,font=\normalfont\bfseries\color{CornellReferenceTeal}]
\item[31.11.2.1] \textbf{Overview.} Frames the terminology, applicability, and relationships needed to interpret Overview; detailed rules remain in the following subclauses.
\item[31.11.2.2] \textbf{Construction and destruction.} Specifies the facility family Construction and destruction. Read its declarations together with mandates, preconditions, effects, results, exceptions, complexity, and lifetime rules.
\item[31.11.2.3] \textbf{Assignment and swap.} Governs state transfer or exchange for Assignment and swap; check self-operation, validity after movement, exception behavior, and invalidation.
\item[31.11.2.4] \textbf{Member functions.} Specifies the facility family Member functions. Read its declarations together with mandates, preconditions, effects, results, exceptions, complexity, and lifetime rules.
\item[31.11.2.5] \textbf{Overridden virtual functions.} Specifies the facility family Overridden virtual functions. Read its declarations together with mandates, preconditions, effects, results, exceptions, complexity, and lifetime rules.
\item[31.11.2.6] \textbf{Specialized algorithms.} Defines the observable result of Specialized algorithms together with applicable iterator-domain, ordering, complexity, invalidation, and exception conditions.
\item[31.11.3.1] \textbf{Overview.} Frames the terminology, applicability, and relationships needed to interpret Overview; detailed rules remain in the following subclauses.
\item[31.11.3.2] \textbf{Construction and destruction.} Specifies the facility family Construction and destruction. Read its declarations together with mandates, preconditions, effects, results, exceptions, complexity, and lifetime rules.
\item[31.11.3.3] \textbf{Member functions.} Specifies the facility family Member functions. Read its declarations together with mandates, preconditions, effects, results, exceptions, complexity, and lifetime rules.
\end{description}
\Needspace{8\baselineskip}
\subsection*{31.12 File systems}
\begin{longtable}{@{}>{\RaggedRight\arraybackslash\columncolor{CornellReferenceCueGray}}p{0.28\textwidth}>{\RaggedRight\arraybackslash}p{0.64\textwidth}@{}}
\rowcolor{CornellReferenceNavy}\color{white}\textbf{Cue / recall prompt} & \color{white}\textbf{Notes}\\\endfirsthead
\rowcolor{CornellReferenceNavy}\color{white}\textbf{Cue / recall prompt} & \color{white}\textbf{Notes (continued)}\\\endhead
\arrayrulecolor{CornellReferenceLineGray}
\textbf{What contract governs General?}\\[-1mm]{\scriptsize\CornellClauseRef{31.12.1}} & Frames the terminology, applicability, and relationships needed to interpret General; detailed rules remain in the following subclauses.\\\hline
\textbf{What contract governs Conformance?}\\[-1mm]{\scriptsize\CornellClauseRef{31.12.2}} & Specifies the facility family Conformance. Read its declarations together with mandates, preconditions, effects, results, exceptions, complexity, and lifetime rules.\\\hline
\textbf{Which obligations govern Requirements?}\\[-1mm]{\scriptsize\CornellClauseRef{31.12.3}} & States the admissibility and semantic obligations for Requirements. Separate compile-time participation rules from caller preconditions and runtime effects.\\\hline
\textbf{What contract governs Header <filesystem> synopsis?}\\[-1mm]{\scriptsize\CornellClauseRef{31.12.4}} & Catalogs the declarations made available by Header <filesystem> synopsis. The synopsis is an interface map; the detailed specifications supply constraints and behavior.\\\hline
\textbf{What contract governs Error reporting?}\\[-1mm]{\scriptsize\CornellClauseRef{31.12.5}} & Defines the failure model for Error reporting, including how an error is represented, reported, propagated, or converted into a diagnostic consequence.\\\hline
\textbf{What contract governs Class path?}\\[-1mm]{\scriptsize\CornellClauseRef{31.12.6}} & Defines filesystem semantics for Class path; path format, error reporting, operating-system dependence, iterator state, and races with external changes must be considered.\\\hline
\textbf{What contract governs Class filesystem\_\allowbreak{}error?}\\[-1mm]{\scriptsize\CornellClauseRef{31.12.7}} & Defines the failure model for Class filesystem\_\allowbreak{}error, including how an error is represented, reported, propagated, or converted into a diagnostic consequence.\\\hline
\textbf{What contract governs Enumerations?}\\[-1mm]{\scriptsize\CornellClauseRef{31.12.8}} & Specifies the facility family Enumerations. Read its declarations together with mandates, preconditions, effects, results, exceptions, complexity, and lifetime rules.\\\hline
\textbf{What contract governs Class file\_\allowbreak{}status?}\\[-1mm]{\scriptsize\CornellClauseRef{31.12.9}} & Specifies the facility family Class file\_\allowbreak{}status. Read its declarations together with mandates, preconditions, effects, results, exceptions, complexity, and lifetime rules.\\\hline
\textbf{What contract governs Class directory\_\allowbreak{}entry?}\\[-1mm]{\scriptsize\CornellClauseRef{31.12.10}} & Defines filesystem semantics for Class directory\_\allowbreak{}entry; path format, error reporting, operating-system dependence, iterator state, and races with external changes must be considered.\\\hline
\textbf{What contract governs Class directory\_\allowbreak{}iterator?}\\[-1mm]{\scriptsize\CornellClauseRef{31.12.11}} & Explains the traversal or view contract for Class directory\_\allowbreak{}iterator; prove domain validity, lifetime, sentinel compatibility, and any borrowing or invalidation property.\\\hline
\textbf{What contract governs Class recursive\_\allowbreak{}directory\_\allowbreak{}iterator?}\\[-1mm]{\scriptsize\CornellClauseRef{31.12.12}} & Explains the traversal or view contract for Class recursive\_\allowbreak{}directory\_\allowbreak{}iterator; prove domain validity, lifetime, sentinel compatibility, and any borrowing or invalidation property.\\\hline
\textbf{What contract governs Filesystem operation functions?}\\[-1mm]{\scriptsize\CornellClauseRef{31.12.13}} & Defines the observable result of Filesystem operation functions together with applicable iterator-domain, ordering, complexity, invalidation, and exception conditions.\\\hline
\end{longtable}
\begin{CornellReferenceSummaryBox}File systems supplies the library semantics portion of this clause. The key review move is to connect its local wording to the clause-wide model, then classify each condition by normative force before relying on it.\end{CornellReferenceSummaryBox}
\Needspace{5\baselineskip}\paragraph{Detailed coverage ledger.}
\begin{description}[style=nextline,leftmargin=2.8cm,font=\normalfont\bfseries\color{CornellReferenceTeal}]
\item[31.12.2.1] \textbf{General.} Frames the terminology, applicability, and relationships needed to interpret General; detailed rules remain in the following subclauses.
\item[31.12.2.2] \textbf{POSIX conformance.} Specifies the facility family POSIX conformance. Read its declarations together with mandates, preconditions, effects, results, exceptions, complexity, and lifetime rules.
\item[31.12.2.3] \textbf{Operating system dependent behavior conformance.} Specifies the facility family Operating system dependent behavior conformance. Read its declarations together with mandates, preconditions, effects, results, exceptions, complexity, and lifetime rules.
\item[31.12.2.4] \textbf{File system race behavior.} Specifies the facility family File system race behavior. Read its declarations together with mandates, preconditions, effects, results, exceptions, complexity, and lifetime rules.
\item[31.12.6.1] \textbf{General.} Frames the terminology, applicability, and relationships needed to interpret General; detailed rules remain in the following subclauses.
\item[31.12.6.2] \textbf{Generic pathname format.} Specifies the facility family Generic pathname format. Read its declarations together with mandates, preconditions, effects, results, exceptions, complexity, and lifetime rules.
\item[31.12.6.3] \textbf{Conversions.} Specifies source and destination conditions for Conversions, the resulting type or value category, and any information-loss, pointer, qualification, or lifetime consequence.
\item[31.12.6.3.1] \textbf{Argument format conversions.} Specifies source and destination conditions for Argument format conversions, the resulting type or value category, and any information-loss, pointer, qualification, or lifetime consequence.
\item[31.12.6.3.2] \textbf{Type and encoding conversions.} Specifies source and destination conditions for Type and encoding conversions, the resulting type or value category, and any information-loss, pointer, qualification, or lifetime consequence.
\item[31.12.6.4] \textbf{Requirements.} States the admissibility and semantic obligations for Requirements. Separate compile-time participation rules from caller preconditions and runtime effects.
\item[31.12.6.5] \textbf{Members.} Specifies the facility family Members. Read its declarations together with mandates, preconditions, effects, results, exceptions, complexity, and lifetime rules.
\item[31.12.6.5.1] \textbf{Constructors.} Specifies how Constructors establishes object state, including participation constraints, initialization effects, and exceptional exits.
\item[31.12.6.5.2] \textbf{Assignments.} Governs state transfer or exchange for Assignments; check self-operation, validity after movement, exception behavior, and invalidation.
\item[31.12.6.5.3] \textbf{Appends.} Specifies the facility family Appends. Read its declarations together with mandates, preconditions, effects, results, exceptions, complexity, and lifetime rules.
\item[31.12.6.5.4] \textbf{Concatenation.} Specifies the facility family Concatenation. Read its declarations together with mandates, preconditions, effects, results, exceptions, complexity, and lifetime rules.
\item[31.12.6.5.5] \textbf{Modifiers.} Specifies the facility family Modifiers. Read its declarations together with mandates, preconditions, effects, results, exceptions, complexity, and lifetime rules.
\item[31.12.6.5.6] \textbf{Native format observers.} Specifies the facility family Native format observers. Read its declarations together with mandates, preconditions, effects, results, exceptions, complexity, and lifetime rules.
\item[31.12.6.5.7] \textbf{Generic format observers.} Specifies the facility family Generic format observers. Read its declarations together with mandates, preconditions, effects, results, exceptions, complexity, and lifetime rules.
\item[31.12.6.5.8] \textbf{Compare.} Specifies the facility family Compare. Read its declarations together with mandates, preconditions, effects, results, exceptions, complexity, and lifetime rules.
\item[31.12.6.5.9] \textbf{Decomposition.} Specifies the facility family Decomposition. Read its declarations together with mandates, preconditions, effects, results, exceptions, complexity, and lifetime rules.
\item[31.12.6.5.10] \textbf{Query.} Specifies the facility family Query. Read its declarations together with mandates, preconditions, effects, results, exceptions, complexity, and lifetime rules.
\item[31.12.6.5.11] \textbf{Generation.} Specifies the facility family Generation. Read its declarations together with mandates, preconditions, effects, results, exceptions, complexity, and lifetime rules.
\item[31.12.6.6] \textbf{Iterators.} Explains the traversal or view contract for Iterators; prove domain validity, lifetime, sentinel compatibility, and any borrowing or invalidation property.
\item[31.12.6.7] \textbf{Inserter and extractor.} Specifies the facility family Inserter and extractor. Read its declarations together with mandates, preconditions, effects, results, exceptions, complexity, and lifetime rules.
\item[31.12.6.8] \textbf{Non-member functions.} Specifies the facility family Non-member functions. Read its declarations together with mandates, preconditions, effects, results, exceptions, complexity, and lifetime rules.
\item[31.12.6.9] \textbf{Hash support.} Specifies the facility family Hash support. Read its declarations together with mandates, preconditions, effects, results, exceptions, complexity, and lifetime rules.
\item[31.12.7.1] \textbf{General.} Frames the terminology, applicability, and relationships needed to interpret General; detailed rules remain in the following subclauses.
\item[31.12.7.2] \textbf{Members.} Specifies the facility family Members. Read its declarations together with mandates, preconditions, effects, results, exceptions, complexity, and lifetime rules.
\item[31.12.8.1] \textbf{Enum path::format.} Specifies the facility family Enum path::format. Read its declarations together with mandates, preconditions, effects, results, exceptions, complexity, and lifetime rules.
\item[31.12.8.2] \textbf{Enum class file\_\allowbreak{}type.} Specifies the facility family Enum class file\_\allowbreak{}type. Read its declarations together with mandates, preconditions, effects, results, exceptions, complexity, and lifetime rules.
\item[31.12.8.3] \textbf{Enum class copy\_\allowbreak{}options.} Specifies the facility family Enum class copy\_\allowbreak{}options. Read its declarations together with mandates, preconditions, effects, results, exceptions, complexity, and lifetime rules.
\item[31.12.8.4] \textbf{Enum class perms.} Specifies the facility family Enum class perms. Read its declarations together with mandates, preconditions, effects, results, exceptions, complexity, and lifetime rules.
\item[31.12.8.5] \textbf{Enum class perm\_\allowbreak{}options.} Specifies the facility family Enum class perm\_\allowbreak{}options. Read its declarations together with mandates, preconditions, effects, results, exceptions, complexity, and lifetime rules.
\item[31.12.8.6] \textbf{Enum class directory\_\allowbreak{}options.} Defines filesystem semantics for Enum class directory\_\allowbreak{}options; path format, error reporting, operating-system dependence, iterator state, and races with external changes must be considered.
\item[31.12.9.1] \textbf{General.} Frames the terminology, applicability, and relationships needed to interpret General; detailed rules remain in the following subclauses.
\item[31.12.9.2] \textbf{Constructors.} Specifies how Constructors establishes object state, including participation constraints, initialization effects, and exceptional exits.
\item[31.12.9.3] \textbf{Observers.} Specifies the facility family Observers. Read its declarations together with mandates, preconditions, effects, results, exceptions, complexity, and lifetime rules.
\item[31.12.9.4] \textbf{Modifiers.} Specifies the facility family Modifiers. Read its declarations together with mandates, preconditions, effects, results, exceptions, complexity, and lifetime rules.
\item[31.12.10.1] \textbf{General.} Frames the terminology, applicability, and relationships needed to interpret General; detailed rules remain in the following subclauses.
\item[31.12.10.2] \textbf{Constructors.} Specifies how Constructors establishes object state, including participation constraints, initialization effects, and exceptional exits.
\item[31.12.10.3] \textbf{Modifiers.} Specifies the facility family Modifiers. Read its declarations together with mandates, preconditions, effects, results, exceptions, complexity, and lifetime rules.
\item[31.12.10.4] \textbf{Observers.} Specifies the facility family Observers. Read its declarations together with mandates, preconditions, effects, results, exceptions, complexity, and lifetime rules.
\item[31.12.10.5] \textbf{Inserter.} Specifies the facility family Inserter. Read its declarations together with mandates, preconditions, effects, results, exceptions, complexity, and lifetime rules.
\item[31.12.11.1] \textbf{General.} Frames the terminology, applicability, and relationships needed to interpret General; detailed rules remain in the following subclauses.
\item[31.12.11.2] \textbf{Members.} Specifies the facility family Members. Read its declarations together with mandates, preconditions, effects, results, exceptions, complexity, and lifetime rules.
\item[31.12.11.3] \textbf{Non-member functions.} Specifies the facility family Non-member functions. Read its declarations together with mandates, preconditions, effects, results, exceptions, complexity, and lifetime rules.
\item[31.12.12.1] \textbf{General.} Frames the terminology, applicability, and relationships needed to interpret General; detailed rules remain in the following subclauses.
\item[31.12.12.2] \textbf{Members.} Specifies the facility family Members. Read its declarations together with mandates, preconditions, effects, results, exceptions, complexity, and lifetime rules.
\item[31.12.12.3] \textbf{Non-member functions.} Specifies the facility family Non-member functions. Read its declarations together with mandates, preconditions, effects, results, exceptions, complexity, and lifetime rules.
\item[31.12.13.1] \textbf{General.} Frames the terminology, applicability, and relationships needed to interpret General; detailed rules remain in the following subclauses.
\item[31.12.13.2] \textbf{Absolute.} Specifies the facility family Absolute. Read its declarations together with mandates, preconditions, effects, results, exceptions, complexity, and lifetime rules.
\item[31.12.13.3] \textbf{Canonical.} Specifies the facility family Canonical. Read its declarations together with mandates, preconditions, effects, results, exceptions, complexity, and lifetime rules.
\item[31.12.13.4] \textbf{Copy.} Specifies the facility family Copy. Read its declarations together with mandates, preconditions, effects, results, exceptions, complexity, and lifetime rules.
\item[31.12.13.5] \textbf{Copy file.} Specifies the facility family Copy file. Read its declarations together with mandates, preconditions, effects, results, exceptions, complexity, and lifetime rules.
\item[31.12.13.6] \textbf{Copy symlink.} Specifies the facility family Copy symlink. Read its declarations together with mandates, preconditions, effects, results, exceptions, complexity, and lifetime rules.
\item[31.12.13.7] \textbf{Create directories.} Specifies the facility family Create directories. Read its declarations together with mandates, preconditions, effects, results, exceptions, complexity, and lifetime rules.
\item[31.12.13.8] \textbf{Create directory.} Defines filesystem semantics for Create directory; path format, error reporting, operating-system dependence, iterator state, and races with external changes must be considered.
\item[31.12.13.9] \textbf{Create directory symlink.} Defines filesystem semantics for Create directory symlink; path format, error reporting, operating-system dependence, iterator state, and races with external changes must be considered.
\item[31.12.13.10] \textbf{Create hard link.} Specifies the facility family Create hard link. Read its declarations together with mandates, preconditions, effects, results, exceptions, complexity, and lifetime rules.
\item[31.12.13.11] \textbf{Create symlink.} Specifies the facility family Create symlink. Read its declarations together with mandates, preconditions, effects, results, exceptions, complexity, and lifetime rules.
\item[31.12.13.12] \textbf{Current path.} Specifies the facility family Current path. Read its declarations together with mandates, preconditions, effects, results, exceptions, complexity, and lifetime rules.
\item[31.12.13.13] \textbf{Equivalent.} Specifies the facility family Equivalent. Read its declarations together with mandates, preconditions, effects, results, exceptions, complexity, and lifetime rules.
\item[31.12.13.14] \textbf{Exists.} Specifies the facility family Exists. Read its declarations together with mandates, preconditions, effects, results, exceptions, complexity, and lifetime rules.
\item[31.12.13.15] \textbf{File size.} Specifies the facility family File size. Read its declarations together with mandates, preconditions, effects, results, exceptions, complexity, and lifetime rules.
\item[31.12.13.16] \textbf{Hard link count.} Specifies the facility family Hard link count. Read its declarations together with mandates, preconditions, effects, results, exceptions, complexity, and lifetime rules.
\item[31.12.13.17] \textbf{Is block file.} Defines the blocking and synchronization contract for Is block file; ownership, wake conditions, timeout behavior, and happens-before edges must be analyzed together.
\item[31.12.13.18] \textbf{Is character file.} Specifies text-sequence behavior for Is character file; establish character type, extent, ownership, null termination, encoding assumptions, and invalidation before use.
\item[31.12.13.19] \textbf{Is directory.} Defines filesystem semantics for Is directory; path format, error reporting, operating-system dependence, iterator state, and races with external changes must be considered.
\item[31.12.13.20] \textbf{Is empty.} Specifies the facility family Is empty. Read its declarations together with mandates, preconditions, effects, results, exceptions, complexity, and lifetime rules.
\item[31.12.13.21] \textbf{Is fifo.} Specifies the facility family Is fifo. Read its declarations together with mandates, preconditions, effects, results, exceptions, complexity, and lifetime rules.
\item[31.12.13.22] \textbf{Is other.} Specifies the facility family Is other. Read its declarations together with mandates, preconditions, effects, results, exceptions, complexity, and lifetime rules.
\item[31.12.13.23] \textbf{Is regular file.} Specifies the facility family Is regular file. Read its declarations together with mandates, preconditions, effects, results, exceptions, complexity, and lifetime rules.
\item[31.12.13.24] \textbf{Is socket.} Specifies the facility family Is socket. Read its declarations together with mandates, preconditions, effects, results, exceptions, complexity, and lifetime rules.
\item[31.12.13.25] \textbf{Is symlink.} Specifies the facility family Is symlink. Read its declarations together with mandates, preconditions, effects, results, exceptions, complexity, and lifetime rules.
\item[31.12.13.26] \textbf{Last write time.} Specifies the facility family Last write time. Read its declarations together with mandates, preconditions, effects, results, exceptions, complexity, and lifetime rules.
\item[31.12.13.27] \textbf{Permissions.} Specifies the facility family Permissions. Read its declarations together with mandates, preconditions, effects, results, exceptions, complexity, and lifetime rules.
\item[31.12.13.28] \textbf{Proximate.} Specifies the facility family Proximate. Read its declarations together with mandates, preconditions, effects, results, exceptions, complexity, and lifetime rules.
\item[31.12.13.29] \textbf{Read symlink.} Specifies the facility family Read symlink. Read its declarations together with mandates, preconditions, effects, results, exceptions, complexity, and lifetime rules.
\item[31.12.13.30] \textbf{Relative.} Specifies the facility family Relative. Read its declarations together with mandates, preconditions, effects, results, exceptions, complexity, and lifetime rules.
\item[31.12.13.31] \textbf{Remove.} Specifies the facility family Remove. Read its declarations together with mandates, preconditions, effects, results, exceptions, complexity, and lifetime rules.
\item[31.12.13.32] \textbf{Remove all.} Specifies the facility family Remove all. Read its declarations together with mandates, preconditions, effects, results, exceptions, complexity, and lifetime rules.
\item[31.12.13.33] \textbf{Rename.} Specifies the facility family Rename. Read its declarations together with mandates, preconditions, effects, results, exceptions, complexity, and lifetime rules.
\item[31.12.13.34] \textbf{Resize file.} Specifies the facility family Resize file. Read its declarations together with mandates, preconditions, effects, results, exceptions, complexity, and lifetime rules.
\item[31.12.13.35] \textbf{Space.} Specifies the facility family Space. Read its declarations together with mandates, preconditions, effects, results, exceptions, complexity, and lifetime rules.
\item[31.12.13.36] \textbf{Status.} Specifies the facility family Status. Read its declarations together with mandates, preconditions, effects, results, exceptions, complexity, and lifetime rules.
\item[31.12.13.37] \textbf{Status known.} Specifies the facility family Status known. Read its declarations together with mandates, preconditions, effects, results, exceptions, complexity, and lifetime rules.
\item[31.12.13.38] \textbf{Symlink status.} Specifies the facility family Symlink status. Read its declarations together with mandates, preconditions, effects, results, exceptions, complexity, and lifetime rules.
\item[31.12.13.39] \textbf{Temporary directory path.} Defines filesystem semantics for Temporary directory path; path format, error reporting, operating-system dependence, iterator state, and races with external changes must be considered.
\item[31.12.13.40] \textbf{Weakly canonical.} Specifies the facility family Weakly canonical. Read its declarations together with mandates, preconditions, effects, results, exceptions, complexity, and lifetime rules.
\end{description}
\Needspace{8\baselineskip}
\subsection*{31.13 C library files}
\begin{longtable}{@{}>{\RaggedRight\arraybackslash\columncolor{CornellReferenceCueGray}}p{0.28\textwidth}>{\RaggedRight\arraybackslash}p{0.64\textwidth}@{}}
\rowcolor{CornellReferenceNavy}\color{white}\textbf{Cue / recall prompt} & \color{white}\textbf{Notes}\\\endfirsthead
\rowcolor{CornellReferenceNavy}\color{white}\textbf{Cue / recall prompt} & \color{white}\textbf{Notes (continued)}\\\endhead
\arrayrulecolor{CornellReferenceLineGray}
\textbf{What contract governs Header <cstdio> synopsis?}\\[-1mm]{\scriptsize\CornellClauseRef{31.13.1}} & Catalogs the declarations made available by Header <cstdio> synopsis. The synopsis is an interface map; the detailed specifications supply constraints and behavior.\\\hline
\textbf{What contract governs Header <cinttypes> synopsis?}\\[-1mm]{\scriptsize\CornellClauseRef{31.13.2}} & Catalogs the declarations made available by Header <cinttypes> synopsis. The synopsis is an interface map; the detailed specifications supply constraints and behavior.\\\hline
\end{longtable}
\begin{CornellReferenceSummaryBox}C library files supplies the library semantics portion of this clause. The key review move is to connect its local wording to the clause-wide model, then classify each condition by normative force before relying on it.\end{CornellReferenceSummaryBox}
\section{Language-Lawyer and Library-Usage Pitfalls}
\begin{CornellReferencePitfallBox}\begin{itemize}[label=\textcolor{CornellReferenceAmber}{$\blacktriangleright$}]
\item Reading a synopsis as the complete behavioral contract.
\item Ignoring mandates, preconditions, exception behavior, or complexity requirements.
\item Using an exposition-only name as though it were public API.
\item Assuming invalidation, ownership, or thread-safety properties that are not stated.
\item Treating successful compilation as proof that semantic requirements and preconditions are met.
\end{itemize}\end{CornellReferencePitfallBox}
\section{Programmer and Implementation Checklist}
\begin{CornellReferenceChecklistBox}\begin{itemize}[label=$\square$]
\item Identify the declaring header or module exposure for each facility.
\item Check template arguments and mandates before evaluating an operation.
\item Establish every precondition at the call site.
\item Record effects, returned value, postconditions, and exceptions separately.
\item Audit iterator, reference, pointer, and object lifetime after mutation.
\item Verify stated complexity and synchronization assumptions.
\item For \CornellClauseRef{31.1}, verify general against the precise rule category and all cited dependencies.
\item For \CornellClauseRef{31.2}, verify iostreams requirements against the precise rule category and all cited dependencies.
\item For \CornellClauseRef{31.3}, verify forward declarations against the precise rule category and all cited dependencies.
\item For \CornellClauseRef{31.4}, verify standard iostream objects against the precise rule category and all cited dependencies.
\end{itemize}\end{CornellReferenceChecklistBox}
\section{Clause Synthesis}
\begin{CornellReferenceOrientationBox}Defines stream state, formatting, buffering, stream buffers, string and span streams, file streams, synchronized output, filesystem paths and operations, and C-library file integration. The formatted and unformatted I/O model plus portable filesystem access. A sound review distinguishes syntax and participation from runtime preconditions, keeps notes and examples non-mandatory, and follows cross-clause dependencies before making a portability claim.\end{CornellReferenceOrientationBox}
\section{Self-Test Questions}
\begin{enumerate}
\item What role does General play in the clause's overall model?
\item Which normative distinctions must be kept separate when applying Iostreams requirements?
\item How would you audit a program or implementation that depends on Forward declarations?
\item Compare Standard iostream objects with Iostreams base classes; what different question does each answer?
\item What role does Iostreams base classes play in the clause's overall model?
\item Which normative distinctions must be kept separate when applying Stream buffers?
\item How would you audit a program or implementation that depends on Formatting and manipulators?
\item Compare String-based streams with Span-based streams; what different question does each answer?
\item What role does Span-based streams play in the clause's overall model?
\item Which normative distinctions must be kept separate when applying File-based streams?
\item How would you audit a program or implementation that depends on Synchronized output streams?
\item Compare File systems with C library files; what different question does each answer?
\item What role does C library files play in the clause's overall model?
\item Which normative distinctions must be kept separate when applying Imbue limitations?
\item Scenario: a program relies on an unstated implementation behavior while using Types. What must be checked before calling the program portable?
\item Scenario: a program relies on an unstated implementation behavior while using Positioning type limitations. What must be checked before calling the program portable?
\end{enumerate}
\clearpage\section{Answer Key}
\begin{enumerate}
\item General is one part of the clause's library semantics model. Its role is understood by connecting the local rule in 31.12.13.1 to the clause orientation and its dependent concepts.
\item Keep formation and well-formedness separate from runtime preconditions and effects. Also distinguish mandatory wording from notes, implementation choice from undefined behavior, and interface declarations from detailed contracts; see 31.2.
\item Audit Forward declarations by locating the controlling wording in 31.3, listing inputs and type constraints, proving any preconditions, recording effects and failure behavior, and checking lifetime, invalidation, complexity, and synchronization where applicable.
\item Standard iostream objects and Iostreams base classes occupy different steps or facility groups. Analyze each under its own subclause, then follow the explicit dependency between them rather than transferring one set of guarantees to the other; begin at 31.4.
\item Iostreams base classes is one part of the clause's library semantics model. Its role is understood by connecting the local rule in 31.5 to the clause orientation and its dependent concepts.
\item Keep formation and well-formedness separate from runtime preconditions and effects. Also distinguish mandatory wording from notes, implementation choice from undefined behavior, and interface declarations from detailed contracts; see 31.6.
\item Audit Formatting and manipulators by locating the controlling wording in 31.7, listing inputs and type constraints, proving any preconditions, recording effects and failure behavior, and checking lifetime, invalidation, complexity, and synchronization where applicable.
\item String-based streams and Span-based streams occupy different steps or facility groups. Analyze each under its own subclause, then follow the explicit dependency between them rather than transferring one set of guarantees to the other; begin at 31.8.
\item Span-based streams is one part of the clause's library semantics model. Its role is understood by connecting the local rule in 31.9 to the clause orientation and its dependent concepts.
\item Keep formation and well-formedness separate from runtime preconditions and effects. Also distinguish mandatory wording from notes, implementation choice from undefined behavior, and interface declarations from detailed contracts; see 31.10.
\item Audit Synchronized output streams by locating the controlling wording in 31.11, listing inputs and type constraints, proving any preconditions, recording effects and failure behavior, and checking lifetime, invalidation, complexity, and synchronization where applicable.
\item File systems and C library files occupy different steps or facility groups. Analyze each under its own subclause, then follow the explicit dependency between them rather than transferring one set of guarantees to the other; begin at 31.12.
\item C library files is one part of the clause's library semantics model. Its role is understood by connecting the local rule in 31.13 to the clause orientation and its dependent concepts.
\item Keep formation and well-formedness separate from runtime preconditions and effects. Also distinguish mandatory wording from notes, implementation choice from undefined behavior, and interface declarations from detailed contracts; see 31.2.1.
\item First identify the exact requirement in 31.5.2.2, then classify the relied-on behavior as required, unspecified, implementation-defined, conditionally supported, or undefined. Portability requires satisfying the program obligations and avoiding dependence on a choice not guaranteed in the target environments.
\item First identify the exact requirement in 31.2.3, then classify the relied-on behavior as required, unspecified, implementation-defined, conditionally supported, or undefined. Portability requires satisfying the program obligations and avoiding dependence on a choice not guaranteed in the target environments.
\end{enumerate}
\section{Key-Term Recap}
\begin{longtable}{@{}>{\RaggedRight\arraybackslash}p{0.28\textwidth}>{\RaggedRight\arraybackslash}p{0.64\textwidth}@{}}\toprule\textbf{Term} & \textbf{Working meaning}\\\midrule\endfirsthead\toprule\textbf{Term} & \textbf{Working meaning (continued)}\\\midrule\endhead\bottomrule\endfoot
ios\_\allowbreak{}base & The shared formatting and state base for stream classes.\\
stream buffer & The controlled character-sequence interface underlying streams.\\
sentry & A helper that performs an operation's preparatory and cleanup work.\\
iostate & The bitmask describing end-of-file, failure, bad state, or good state.\\
syncstream & A buffered wrapper that emits character groups without interleaving.\\
filesystem path & A pathname abstraction with native and generic formats.\\
filesystem\_\allowbreak{}error & An exception carrying filesystem error and optional path information.\\
\end{longtable}
\CornellTakeaway{Reliable I/O separates stream state, formatting, buffering, device interaction, encoding, and filesystem race behavior.}
\end{document}
